%% kapa: exam=1 round=21 year=2010 period=1 form=A subject=civil_law subject_ko=민법 seq=30 number=30
%% source: Q-Net 감정평가사 기출문제 (cst00309) artlSeq=5024181 file=extracted/1008/5024181/2010년 제21회 감정평가사 1차시험 기출문제/1교시.hwp
%% conversion: kordoc 4.24.51 -> tools/gen1.py (auto); review=proofread-sonnet; flags=none
\begin{question}{30}
甲은 乙로부터 500만원을 빌리면서 그 채무를 담보하기 위해서 자신의 반지에 질권을 설정하려고 한다. 다음 설명으로 옳은 것은? (다툼이 있으면 판례에 의함)
\begin{choices}
  \item 甲은 자신이 그 반지를 직접점유하면서 乙이 질권을 취득하게 할 수 있다.
  \item 반지가 甲의 소유가 아니라 甲이 丙으로부터 빌린 것인 경우에도, 乙이 선의ㆍ무과실로 반지를 인도받았다면 질권을 선의취득한다.
  \item 甲이 乙과의 질권설정계약 이전에 제3자 丁과 “반지를 타인에게 담보로 제공하지 않는다.”는 약정을 하였다면, 乙은 질권을 취득하지 못한다.
  \item 변제기가 도래한 후에 “변제에 갈음하여 乙이 반지의 소유권을 취득하기로 한다.”고 한 甲ㆍ乙간의 약정은 무효이다.
  \item 질권 설정 후 제3자가 반지를 침탈한 경우, 乙은 질권 자체에 기한 물권적 청구권을 행사할 수 있을 뿐 점유회수청구권을 행사할 수는 없다.
\end{choices}
\end{question}
\answer{2}
